\section{Discussion}\label{sec:discussion}

\subsection{What the separation is, and what it is not}\label{sec:scope}

The scalarized operators $\{E_w\}_{w\in\Wp}$ formalise one compensatory certification rule: a plan is assessed through a weighted aggregate of the declared normalized margins, with the stated disturbances respected. The weak-sustainability literature is broader than this operator, encompassing aggregate wealth, nondeclining comprehensive consumption, shadow prices and intertemporal investment conditions \citep{Solow1974,Asheim1994,DasguptaMaler2000,Neumayer2013}. Strong-sustainability frameworks are also broader than the rule that every coordinate remain nonnegative at every instant; critical-natural-capital and safe-operating-space approaches include thresholds, irreversibility and resilience \citep{Ekins2003,Rockstrom2009,Raworth2012,DoyenGajardo2020}. The two operators are therefore idealized certification protocols, not definitions of the full doctrines.

The established gap is one of plan selection: every weighting may license some plan even though no single plan satisfies all typed floors. The quantifier order has an information interpretation only if the plan can actually be chosen after the assessment weight is known; the mathematical result itself does not assume such information or describe auditor behavior. The same mechanism can occur for static alternative vectors. In the present application, paths and disturbances add questions about transient breaches, terminal conditions, and which control extensions are implementable. The finite-menu theorem's convex geometry is a menu-selection result under margin-first scalarization; it is not evidence that the mechanism is inherently dynamic.

The rational witness is relative to its four-plan menu and declared disturbance class. A larger or differently implemented menu can shrink the gap. The reserve is deliberately excluded from the aggregate, so adding it to a real indicator could change the result. The action-indexed blend closes the gap only under the stated linear-flow assumption; the shared-shock blend improves the worst margin at the canonical point from $-4/5$ to $-1/20$ but remains infeasible. The complete-excursion result is conditional on a union-of-full-tubes policy class, while a distinct three-switch no-shock schedule is safe at the canonical margins. No general claim about finite-rate switching follows. The coupled-shock results are reading- and depth-conditional (SI~S4). Finally, all formal results concern finite-horizon exact-tube data; infinite horizons, stochastic constraints, partial observation and endogenous event times are outside the main model.

\subsection{Acceptance semantics in multi-criteria decision aiding}\label{sec:mcda}

The protocols in Table~\ref{tab:protocols} are sorting rules: they specify when a transition is classified as accepted or rejected, the sorting problem of decision aiding \citep{Roy1996,Vincke1992}. The licensing thresholds of Theorem~\ref{thm:witness} are weight-sensitivity intervals, analogous to sensitivity analyses for composite-indicator rankings \citep{Nardo2008,Becker2017}, but the object is which plan is licensed rather than a rank. Existing MCDA work analyzes compensability and outranking at fixed profiles \citep{Cinelli2014,Schar2025}; here the fixed-profile full-cone result is lossless (Lemma~\ref{lem:cone}), and the additional gap arises when acceptance is quantified over a finite menu.

The here-and-now/wait-and-see distinction in adjustable robust optimization is a useful analogy for common-plan versus weight-adaptive acceptance \citep{BenTal2004}. It is not a literal uncertainty model: the weight is an assessment parameter, and the result does not reverse the usual set inclusion between static and adaptive policy classes on a common model. The companion biocapacity study and obstruction-calculus paper address related applications and substitution certificates \citep{Abaee2026ecomod,Abaee2026obs}; they are not premises of the present proofs.

\subsection{Design tests for composite indicators}\label{sec:design}

For indicator developers, the results suggest four checks. Each follows from a stated theorem and applies only within its declared model.
\begin{enumerate}[leftmargin=2em]
\item \textbf{Fixed-profile losslessness.} Does the declared weight family detect every component violation on a fixed trajectory? The full nonnegative cone does (Lemma~\ref{lem:cone}); a restricted nonempty family requires its own characterization (Proposition~\ref{prop:monotone}).
\item \textbf{Protocol quantifiers.} Does certification require one plan to work for every weighting, or permit plan reselection by weight? If the latter, report the plan ranges and the menu-relative gap.
\item \textbf{Path and endpoint checks.} Is the assessment evaluated on the full tube or only at sampled endpoints? The witness's fast plan is typed-endpoint-admissible even where its tube breaches a floor, so endpoint-only review can miss a transient violation.
\item \textbf{Menu and control implementation.} Is the menu exhaustive, and can fractional controls actually be implemented? The convex-hull theorem is margin-first; the physical blend requires a separately verified control law. Declare the disturbance class and do not infer feasibility from an abstract convex combination alone.
\end{enumerate}

Where floors are separately binding, reporting per-floor path minima, the reserve and the acceptance protocol allows this class of discrepancy to be reconstructed. Floors alone do not determine whether a reserve-financed rescue is available, and an aggregate without path and reserve data does not recover the typed verdict on the witness family. This is a conditional design recommendation, not a universal prescription for sustainability assessment. A ten-point checklist, a minimum report and empirical-data requirements appear in SI~S6.

\subsection{Limitations}\label{sec:limitations}

\begin{enumerate}[label=(\roman*),leftmargin=2em]
\item The separation is an existence result on one declared finite menu; it does not estimate the size or frequency of a gap in applications.
\item The finite-menu convex-hull theorem assumes margin-first requirements $m_a(w)=w\cdot d_a$. The synchronization lemma states when that formula agrees with an exact-tube scalarization.
\item A physically implemented convexified control is an additional model and institutional assumption. Without it, convexification in required-margin space is only geometric.
\item The hold-prefix persistence result is proved for its specified extension. A separate repair-stage example shows one erasure mechanism; it does not imply that arbitrary continuations, coincident endpoint sets, or action coupling erase pathwise gaps.
\item The biomass illustration is synthetic. Its optional mapping to the historical DFO limit implies very large model capacities, and its nonlinear recovery does not return to the abstract successor without additional terminal control. The income-margin paths are stipulated rather than generated from an economic model.
\item CES results use the explicit collapse convention and the normalization $\lambda_i=1+s_i$. The diagonal thresholds are analytic; the off-diagonal formula is stated separately, while no global fixed-sum stationary-point or phase theorem is claimed. The common-shock statements apply only to the specified decomposition.
\item The finite rational verifiers check enumerated instances, not continua. The Lean source formalizes selected components and does not cover the full real-exponent, continuity, or topological steps; version and artifact limitations are detailed in SI~S7.
\item Novelty statements are bounded by the literature comparisons in Section~\ref{sec:related}; no claim is made that composite-index or sustainability scholarship lacks formal analysis.
\end{enumerate}

\section{Conclusions}\label{sec:conclusions}

At a fixed trajectory, aggregation over the full nonnegative cone is lossless. A separation can nevertheless arise when the assessment rule permits each weight to select a different plan. For finite menus under the explicit margin-first hypothesis, the accepted-set gap is the part of the up-closed convex hull of required margins not covered by any single plan's orthant. A rational transition witness realizes a nonempty gap and an exact reserve threshold. This is a menu-selection result that can be static; pathwise tubes and disturbances specify its application to transition certification.

Under the witness's linear-flow assumption, a fractional physical blend closes the action-indexed gap. A complete-excursion concatenation preserves primitive floor violations only under its stipulated union-of-tubes semantics, and another finite-switch extension can be safe. Under the specified shared-shock class, the gap survives an exact $1/2$-cut and the balanced blend improves but does not restore feasibility at the canonical point. On the witness-defined diagonal, every positive-elasticity CES member falsely certifies an exact interval under the stated collapse convention, whereas the Leontief limit reproduces typed acceptance over the full weight family. These are model-specific statements, not general claims about randomization, all switching policies, or real fishery dynamics.

For transition reporting, the practical conclusion is to declare the menu, the disturbance class, the normalized margin scales, the plan-selection rule, the path minima, the terminal treatment, and any resource needed for rescue. An aggregate can remain positive while a typed floor is breached, but the example does not show that every composite index is unsound or that a particular policy doctrine is normatively required. The benchmark is a transparent dimensional illustration anchored to a historical DFO value, not a fitted or predictive cod model.
